
\chapter*{怎样阅读这份讲义}
\addcontentsline{toc}{chapter}{怎样阅读这份讲义}
把几个数替换成它们的平均值，是一个几乎不需要准备知识就能提出的操作。
然而，要求在有限步后\emph{精确}得到相同数值，会把问题带到整数同余、二维矩阵与算术群。
这份讲义围绕一条完整结论展开：对奇素数 $p$，位置数达到 $2p+1$ 以后，
有限可达性由一个最大公因子的素因子完全决定，而 $2p+1$ 不能再降低。

讲义假定读者熟悉线性代数、整除与同余、中国剩余定理，以及群、环和群作用的基本定义。
对本题另外需要的概念，例如素数赋值、局部化环、本原列、射影返回和主同余子群，
都在首次使用时说明。两个较深的算术群定理作为已发表定理引用；
这里完整讲解其假设、结论，以及它们怎样嵌入平均问题的证明。
无需预先学习算术群的一般理论。

\section*{一条主线，两种终止方式}
证明有两种主要工作方式。位置较多时，可以始终选择整数平均，
并保证平方和严格下降。平方和取非负整数值，这就把无限过程排除。
位置接近阈值时，可用于选择的位置太少，必须允许分数。
这时先把配置变为三个等值块，用两个参数描述它，再用一系列真实平均控制这两个参数。
目标是先造出一个和为零的 $p$ 元组，随后由已有的含零收尾完成。

因此，阅读矩阵构造时应始终保留下面的图景：
\[
\text{原始配置}\ \longrightarrow\ \text{三块核心}
\ \longrightarrow\ \text{零和的 }p\text{ 元组}
\ \longrightarrow\ \text{全零配置}.
\]
中间的矩阵只是记录真实操作的工具。它的两列不是额外增加的两个位置，
矩阵的逆也必须另有正向平均序列才能使用。

\section*{各部分怎样衔接}
第1--2讲建立问题和必要条件，并证明阈值的下界。
第3--5讲先独立取得整数下降、$3p$ 个位置和含零状态的解法。
这些是后面构造可以放心调用的终点。
第6讲处理任意输入到三块核心的入口。
第7--9讲集中建立矩阵和同余工具；它们解释后面所有参数构造的共同目的。

第10--14讲按 $n$ 与 $p$ 的关系给出真实返回。
此处每一张计数表都回答一个具体问题：某种矩阵作用能否用现有数量的位置完成。
第15讲补齐宽维数中整数下降的少量容量不足情形。
第16--17讲处理 $p=3,5,7$ 所必需的基例；这些例子属于总定理的证明。
第18讲将所有范围合并。

第一次阅读，建议依次读完第1--9讲，再完整读一遍第10讲的奇数余量构造。
这样便能在一个完整实例中看到“来源分组、返回矩阵、剪切、同余群、终端”的全部环节。
然后阅读其他参数段。若只想查某个具体 $(p,n)$，可先看第18讲的覆盖表，
再回到相应讲次。目录中的命题、公式及表格均有文内引用。

\section*{记号约定}
\begin{longtable}{p{.18\textwidth}p{.74\textwidth}}
\toprule 记号 & 含义\\\midrule
$p$ & 每次平均所选位置数；主体中为奇素数。\\
$n$ & 配置的总位置数，始终保留这些位置。\\
$a^m$ & 在配置记号中表示 $m$ 个值为 $a$ 的位置，不是数值幂。\\
$X$ & 减去全局均值并清分母、除公因子后得到的本原整数配置。\\
$G(X)$ & 全部坐标差的最大公因子。\\
$d_p(n)$ & $n$ 的不同素因子中不等于 $p$ 的个数。\\
$v_\ell(t)$ & 有理数 $t$ 的分子中 $\ell$ 的指数减去分母中的指数。\\
$U(t),L(t)$ & 二阶上、下三角初等剪切矩阵。\\
$R=\Z[1/D]$ & 只允许 $D$ 的素因子出现在分母中的有理数环。\\
$\Gamma(q,R)$ & $\SL_2(R)$ 中模 $qR$ 等于单位矩阵的子群。\\
$\cont(v)$ & 整数列坐标的最大公因子，称为内容。\\
\bottomrule
\end{longtable}
矩阵统一作用在列向量上，复合从右到左。描述平均序列时则按实际时间顺序叙述。
余量统一写成 $r=n-2p$；上半带另外写 $s=n-3p$，于是 $r=p+s$。
各构造中的临时计数变量只在所在小节内使用。

\section*{证明的来源与有限数据}
本讲义将已有完整证明编排为连贯课程，并恢复原稿中已有的分组说明和推导细节。
涉及小参数的有限表按原有结果收入正文及附录，所需方程、状态和生成元一并给出。
有限表的作用限于已明确列出的有限命题；无穷参数的论证由正文中的引理承担。
配套目录中的有限数据供查阅，正文不要求读者先运行程序。
本次编撰的重点是叙述、衔接和排版，沿用已有数学结果。
